\subsection*{A compatible class-three lift, including the rigid tests}

\statement{B.61}{Proved}{The common-image obstruction survives through class three}
Let $G_3=P/\gamma_4P$, let $w$ be B.59's primitive
degree-two vector, and put $D=2545920$.
There are explicitly specified finite words
$A\in H_T$ and $B\in H_5$ whose images in $G_3$ are equal
to a nonidentity element $g$, and such that
\[
 g\bmod\gamma_3P=5091840\,w,
 \qquad
 g\longmapsto1\quad\hbox{in every }
 P\Gamma_{0,S}/\gamma_4(P\Gamma_{0,S}),\quad S\in\mathscr S_7.
                                                               \tag{B.61.1}
\]
The element $g$ also centralizes the image of $A_5$ in
$\Gamma_{0,[12]}/\gamma_4P$.
Their common image therefore generates an infinite cyclic
subgroup in the discrete quotient and an injective copy of
$\widehat{\mathbb Z}$ in $\widehat{G_3}$, contained in both
completed subgroup images and in all twelve class-three
rigid-test kernels.
The two words are distinct in $P$; no nonidentity element
of the actual $\mathcal K$ is obtained.

\begin{proof}
\emph{One common geometric marking.}
Higher-degree calculations must compare the two centralizers
at the same marked configuration.
Let $\varphi=(1+\sqrt5)/2$ and set
\[
 u_0=i\bigl(\varphi+\sqrt{\varphi+2}\bigr),\qquad
 \zeta=\frac{\sqrt5-1}{4}+
                  \frac{i\sqrt{\varphi+2}}2,
 \qquad v_0=-\varphi^2.                                    \tag{B.61.2}
\]
Substitute $u_0$ into the twelve tetrahedral functions of
B.56.  In the cyclic description, put labels $1,3$ at
$0,\infty$, put the first displayed cycle of B.60 at
$1,\zeta,\ldots,\zeta^4$, and the second at
$v_0,v_0\zeta,\ldots,v_0\zeta^4$.
After normalizing labels $2,4,6$ to $\infty,0,1$,
these two labelled configurations agree exactly.
The cyclic normalization is
\[
 z(q)=\frac{q-\zeta^4}{q-1}
             \frac{\zeta-1}{\zeta-\zeta^4},                \tag{B.61.3}
\]
with the usual value at $q=\infty$.
All eleven equalities are checked in the degree-eight field
\[
                  \mathbb Q(i,\sqrt5,\sqrt{\varphi+2}).
\]
The same check verifies $\zeta^5=1$ and the prescribed
two five-cycles.  Together with the given tetrahedral
action, this is the actual marked $A_5$ configuration.

The program \nolinkurl{next_B59_image/certify_full11_monodromy.py}
extends the thirteen exact meridian paths of B.56 to all
eleven finite normalized points.  Exact Sturm isolation
certifies all crossings on their $78$ rational segments.
Each resulting braid word projects \emph{literally} to
the corresponding earlier six-strand word, and all $55$
pair linkings agree with B.58's independently calculated
rational-divisor matrix.
The complete polynomial and interval data are in
\nolinkurl{next_B59_image/full11_exact_certificate.json}.

For the cyclic family, rational approximations to the five
roots, each with error less than $10^{-7}$, permit the
same exact rational Sturm computation.
This approximation is accompanied by a homotopy certificate.
Along every chosen meridian and transport segment, the
parameter has absolute value less than seven and its
squared distance from every rational puncture centre is
at least $10^{-4}$.
Thus the error in any moving--fixed point difference is
less than $8\cdot10^{-7}$, strictly below the separation
from an actual collision.  Moving--moving differences
and the three normalizing points remain separated as well.
Consequently the perturbed family and the actual cyclic
family have the same braid monodromy after their specified
base transport.

Rational interval arithmetic, with radical bounds verified
by squaring, also proves that the algebraic-to-rational
base paths lie in a single crossing-free cell for the
projection multiplied by $1+317i/1000$.
The common left-to-right ordering of the finite labels is
\[
                   (4,11,7,9,3,5,1,12,8,10,6).
                                                               \tag{B.61.4}
\]
The transport from the old tetrahedral base has $16$
crossings; the cyclic base transport has $39$.
The complete checks, six cyclic meridian words and both
transports are recorded in
\nolinkurl{next_B59_image/certify_cyclic_and_transport.py},
its actual output, and its exact JSON certificate.
The resulting common-base words are in
\nolinkurl{next_B59_image/common_base_braid_words.json}.
This establishes the actual embeddings used in the
calculation, including their relative marking.

\emph{The exact nilpotent calculation.}
The Artin action uses
$x_i\mapsto x_ix_{i+1}x_i^{-1}$ and
$x_{i+1}\mapsto x_i$ for a positive braid generator.
Integer Magnus expansions through degree four compute
the first three homogeneous terms of its logarithm,
recorded without denominators as $D_1,2D_2,24D_3$.
The programs \nolinkurl{next_B59_image/nilpotent_derivations.py}
and \nolinkurl{next_B59_image/magnus4.py} derive and check
the truncated conjugation and Baker--Campbell--Hausdorff
formulae.  Their checks include braid relations, inverses,
agreement of successive truncations, products, and the
sign of a weight-three group commutator.
The degree-one terms of all nineteen geometric generators
are independently checked against their pair linkings.

The integral infinitesimal braid presentation gives a
degree-three basis by the standard length-three Lyndon
brackets in
$A_{1k},\ldots,A_{k-1,k}$, for each last strand
$k=2,\ldots,11$.
Its rank is
\[
                 \sum_{r=1}^{10}\frac{r^3-r}{3}=990.
                                                               \tag{B.61.5}
\]
This is the degree-three part of the usual free-group
forgetful filtration; the presentation is also recorded
in \citet[Section 3]{DE_KaabiVershinin2012}.
For completeness, the certificate checks coordinates
through the graded Artin derivations themselves.
For a last strand $k$, the coefficients in its image
of words containing that strand exactly once, at the end,
recover the corresponding free-Lie polynomial on the
preceding strands.  Standard Lyndon leading coefficients
are one.  Subtraction therefore extracts its coordinates,
and reconstruction verifies the \emph{entire} derivation.
This proves the injectivity needed in degrees one through
three, rather than assuming faithfulness of a finite
matrix approximation.
The central full twist affects degree one alone; all
geometric words have zero linking on the fixed pair
$(4,6)$, so the calculation represents $P$ itself.

Let $a,b_0$ be B.60's words.  Their degree-two images are
$2w,w$, respectively, and
\[
                 \delta=ab_0^{-2}\bmod\gamma_4P
                  \in\gamma_3P/\gamma_4P.                  \tag{B.61.6}
\]
The actual correction spaces are the images in this
central layer of $H_T\cap\gamma_3P$ and
$H_5\cap\gamma_3P$.
The cyclic degree-one and degree-two maps are injective,
so $H_5\cap\gamma_3P=\gamma_3H_5$; its $70$ standard
ternary brackets give its correction space.
For $H_T$, the search retains all thirteen generator
logarithms, all $78$ binary brackets, and all ternary
brackets.  Its lower-layer matrix has rank $19$ and
nullity $72$.  Its kernel supplies candidate corrections
while retaining the seven source-abelian directions and
the bracket relations.  The existence assertion below
rests on the explicit integral words and their direct
verification; it does not require a comparison of a
rational correction space with an integral subgroup lattice.
In particular, $H_T\cap\gamma_3P$ is not replaced by
$\gamma_3H_T$.

The required linear equations ask for a tetrahedral
correction $u$ and cyclic correction $v$ with
$u+v=\delta$, and also require the rigid-seven projection
of $u$ to equal that of $a$ in degree three.
For the target containing labels $2,4,6$, this last map
is a literal projection onto $70$ of the chosen Hall
coordinates.  The exact rational solution is checked
against all $990+70$ equations.
The mod-$101$ calculation only selects a convenient
pivot block; it is not used to infer rational membership.
All matrices, pivot inverses and the actual rational
identity are retained in the two
\nolinkurl{degree3_actual_correction_certificate.json} and
\nolinkurl{degree3_rigid_obstruction_certificate.json}
files under \nolinkurl{next_B59_image/}.

\emph{Actual integral words, not merely rational membership.}
The file
\nolinkurl{next_B59_image/class3_common_lift_words.json}
specifies two correction words completely, by integer
powers in a fixed order.
The tetrahedral word $U$ has $13$ generator powers,
$78$ commutator powers and $728$ standard Lyndon
weight-three powers.  The cyclic word $V$ has $39$
weight-three powers.
The program
\nolinkurl{next_B59_image/collect_integral_class3_lift.py}
constructs them using $D=2545920$ and verifies their
collection in the \emph{free} nilpotent group of class
three by the exact BCH formula.  Every displayed power
is an integer; rational roots of group elements are
not used.
Their images satisfy
\[
 U\in H_T\cap\gamma_3P,\qquad V\in\gamma_3H_5,
 \qquad [U]+[V]=D\delta,                                  \tag{B.61.7}
\]
and the rigid degree-three image of $U$ is $D$ times
that of $a$.

Define the promised finite words by
\[
                       A=a^D U^{-1},\qquad
                       B=b_0^{2D}V.                      \tag{B.61.8}
\]
The program
\nolinkurl{next_B59_image/verify_integral_class3_lift.py}
reads the \emph{collected integer word exponents} and
evaluates them again in the actual Artin logarithms.
It checks every tetrahedral generator and commutator
power, all $728$ tetrahedral triple powers, and all
$39$ cyclic triple powers, with full graded reconstruction.
It obtains equality of (B.61.8) through class three and
zero rigid image, independently of the rational-span
calculation.  The actual standard output is retained.
The independent root review reran this word evaluation;
its passing record is \nolinkurl{records/root_review_B61.md}.
Equations (B.61.7)--(B.61.8) prove (B.61.1) for one
rigid-seven set.  Because $A$ is an actual word in
$H_T$, conjugacy by $T$ proves it for all twelve sets,
as in B.56.
The common image centralizes both $T$ and $D$, and hence
the group $\langle T,D\rangle=A_5$.

The common class-two image $2Dw$ is nonzero, so $g$ has
infinite order.  Its profinite power map is injective
because its projection to B.59's completed central lattice
is multiplication by the nonzero integer $2D$ on a
primitive $\widehat{\mathbb Z}$ coordinate.
Continuity gives the asserted common procyclic image
and the twelve completed rigid-test vanishing statements.
Finally $H_T\cap H_5=1$ in $P$.  Equality $A=B$ there
would force their nonzero class-two image to vanish.
Hence they are distinct, with $AB^{-1}\in\gamma_4P$,
and the result makes no assertion about a lift to
$\mathcal K$.
\end{proof}
